\pagestyle{empty}
\tight
\begin{tblr}{width=\textwidth,colspec={|Q[c,m,1.9cm]|X[l,m]|},hlines,vlines,hspan=minimal,row{1}={bg=headblue},rowsep=1pt,colsep=3pt}
\SetCell{c,m}\hfil\logo\hfil & \textbf{POLITEKNIK NEGERI MALANG}\par\textbf{JURUSAN TEKNOLOGI INFORMASI}\par\textbf{PROGRAM STUDI: D4 TEKNIK INFORMATIKA} \\
\SetCell[c=2]{c} \textbf{RENCANA PEMBELAJARAN SEMESTER (RPS)} & \\
\end{tblr}
\begin{longtblr}{width=\textwidth,colspec={|Q[l,m,1.9cm]|X[0.85,l,m]|X[1.45,l,m]|X[0.75,l,m]|X[1.15,l,m]|},hlines,vlines,hspan=minimal,rowsep=1pt,colsep=2pt,row{1,3}={bg=headgray,font=\bfseries}}
MATA KULIAH & KODE & BOBOT (SKS)/jam & SEMESTER & TGL. PENYUSUNAN \\
Dasar Pemrograman & RTI251006 & 2 SKS/4 jam & 1 & 7 Agustus 2025 \\
OTORISASI & Dosen Pengembang RPS & Koordinator RMK & \SetCell[c=2]{l} Ka PRODI &  \\
 & Mungki Astiningrum, S.T., M.Kom.\par Imam Fahrur Rozi, S.T., M.T.\par Vivi Nur Wijayaningrum, S.Kom., M.Kom.\par Mamluatul Hani'ah, S.Kom., M.Kom. & Imam Fahrur Rozi, S.T., M.T. & \SetCell[c=2]{l}{\vspace{8mm}\par Dr. Ely Setyo Astuti, S.T., M.T.} &  \\
\SetCell[r=11]{l} Capaian Pembelajaran (CP) & \SetCell[c=4]{l,bg=headgray,font=\bfseries} Capaian Pembelajaran Lulusan Program Studi (CPL-Prodi) &  &  &  \\
 & CPL07 & \SetCell[c=3]{l} Mampu menerapkan sistem komputasi cerdas dan pengolahan data untuk mendukung penyelesaian masalah komputasi. &  &  \\
 & CPL09 & \SetCell[c=3]{l} Mampu menjelaskan cara kerja sistem komputer dan menerapkan pendekatan komputasi untuk menyelesaikan masalah. &  &  \\
 & \SetCell[c=4]{l,bg=headgray,font=\bfseries} Capaian Pembelajaran mata kuliah (CPL-MK) &  &  &  \\
 & CPMK0704 & \SetCell[c=3]{l} Mahasiswa mampu menerapkan berbagai bahasa dan paradigma pemrograman untuk menyelesaikan masalah komputasi. &  &  \\
 & CPMK0903 & \SetCell[c=3]{l} Mahasiswa mampu menjelaskan prinsip logika komputasi serta menerapkan teknik dasar pemrograman dan algoritma untuk menganalisis dan menyelesaikan permasalahan. &  &  \\
 & \SetCell[c=4]{l,bg=headgray,font=\bfseries} Kemampuan akhir tiap tahapan belajar (Sub-CPMK) &  &  &  \\
 & SCPMK0903-00601 & \SetCell[c=3]{l} Mahasiswa mampu menjelaskan konsep dasar algoritma, version control, dan Kanban serta menganalisis permasalahan sederhana ke dalam bentuk algoritma [C2, A3]. &  &  \\
 & SCPMK0903-00602 & \SetCell[c=3]{l} Mahasiswa mampu menjelaskan tipe data, variabel, input-output, sequence, dan operator serta menerapkannya dalam penulisan algoritma [C4, A3]. &  &  \\
 & SCPMK0704-00603 & \SetCell[c=3]{l} Mahasiswa mampu menuliskan algoritma pemilihan dan perulangan, baik sederhana maupun bersarang, menggunakan flowchart untuk menyelesaikan studi kasus [C4, A3, P2]. &  &  \\
 & SCPMK0704-00604 & \SetCell[c=3]{l} Mahasiswa mampu menyelesaikan studi kasus menggunakan array 1 dan 2 dimensi serta fungsi iteratif dan rekursif dalam algoritma penyelesaian masalah [C4, A3, P2]. &  &  \\
\SetCell[r=2]{l} Penilaian & \SetCell[c=4]{l} *Nilai CPMK didapat dari agregasi nilai SUB-CPMK yang dibebankan &  &  &  \\
 & \SetCell[c=4]{l}{\tiny\begin{tblr}{width=\linewidth,colspec={|Q[c,m,0.1\linewidth]|Q[c,m,0.16\linewidth]|X[l,m]|X[l,m]|X[l,m]|X[l,m]|X[l,m]|Q[c,m,0.08\linewidth]|},hlines,vlines,hspan=minimal,rowsep=1pt,colsep=0.7pt,row{1,2}={bg=headgray,font=\bfseries}}
Kode CPMK & Kode SCPMK & \SetCell[c=5]{c} Bobot per Bentuk Penilaian & & & & & Total Bobot \\
 & & Tugas & Kuis 1 & UTS & Kuis 2 & UAS & \\
\SetCell[r=2]{c} CPMK0903 & SCPMK0903-00601 & 4\%: konsep algoritma, version control, Kanban & 5\%: konsep algoritma & 6\%: konsep dasar algoritma & & & 15\% \\
 & SCPMK0903-00602 & 2\%: tipe data, operator, sequence & 5\%: tipe data dan operator & 8\%: tipe data, operator, sequence & & & 15\% \\
\SetCell[r=2]{c} CPMK0704 & SCPMK0704-00603 & 8\%: flowchart pemilihan dan perulangan & & 16\%: algoritma pemilihan dan perulangan & 5\%: perulangan bersarang & 10\%: kasus pemilihan dan perulangan & 39\% \\
 & SCPMK0704-00604 & 6\%: array dan fungsi & & & 5\%: array 1 dan 2 dimensi & 20\%: kasus array dan fungsi & 31\% \\
\SetCell[c=7]{r}\textbf{Total Per Penilaian} & & & & & & & \textbf{100\%} \\
\end{tblr}} &  &  &  \\
Diskripsi Singkat Mata Kuliah & \SetCell[c=4]{l} Dasar Pemrograman memberikan pengetahuan dan pemahaman konsep dasar algoritma dan dasar pemrograman sehingga mahasiswa memiliki dasar untuk menyelesaikan permasalahan-permasalahan logika dengan menggunakan flowchart dan pseudocode. &  &  &  \\
Bahan Kajian / Materi Pembelajaran & \SetCell[c=4]{l}\begin{enumerate}\item Konsep algoritma, analisis permasalahan, version control, dan Kanban\item Tipe data, variabel, konstanta, nilai, ekspresi, input-output, sequence, dan operator\item Analisa kasus: pemilihan dan perulangan (sederhana dan bersarang)\item Array 1 dan 2 dimensi, fungsi iteratif dan rekursif\end{enumerate} &  &  &  \\
\SetCell[r=2]{l} Pustaka & \SetCell[c=4]{l} \textbf{Utama:}\begin{enumerate}\item Sebesta, Robert. 2016. Concept of Programming Languages, Edisi Global. Addison Wesley.\item Sestoft, Peter. 2017. Programming Language Concepts. Springer.\item T. Henny Febriana Harumy. 2016. Belajar Dasar Algoritma dan Pemrograman C++. Deepublish.\end{enumerate} &  &  &  \\
 & \SetCell[c=4]{l} \textbf{Pendukung:} Rinaldi Munir. 2015. Algoritma dan Pemrograman. Penerbit Informatika. &  &  &  \\
Media Pembelajaran & \SetCell[c=2]{l} \textbf{Software:} Microsoft Office, Adobe Reader, Sublime Text, JDK. &  & \SetCell[c=2]{l} \textbf{Hardware:} PC/Laptop. &  \\
Nama Dosen Pengampu & \SetCell[c=4]{l}\begin{enumerate}\item Mungki Astiningrum, S.T., M.Kom.\item Imam Fahrur Rozi, S.T., M.T.\item Triana Fatmawati, S.T., M.T.\item Mamluatul Hani'ah, S.Kom., M.Kom.\item Rudy Ariyanto, S.T., M.Cs.\item Noprianto, S.Kom., M.Eng.\end{enumerate} &  &  &  \\
Matakuliah Syarat & \SetCell[c=4]{l} - &  &  &  \\
\end{longtblr}
\begin{longtblr}{width=\textwidth,colspec={|Q[c,m,0.06\textwidth]|X[1.35,l,m]|X[1.08,l,m]|X[1.16,l,m]|X[0.7,l,m]|X[1.24,l,m]|X[1.06,l,m]|X[0.98,l,m]|Q[c,m,0.05\textwidth]|},hlines,vlines,hspan=minimal,rowhead=2,row{1,2}={bg=headgreen,font=\bfseries},rowsep=1pt,colsep=1.5pt}
Mgg.\par Ke & Kemampuan Akhir Yang Direncanakan (Sub-CP-MK) & Bahan kajian (Materi Pembelajaran) & Bentuk dan Metode Pembelajaran & Estimasi Waktu & Pengalaman Belajar Mahasiswa & Kriteria \& Bentuk Penilaian & Indikator Penilaian & Bobot Penilaian (\%) \\
(1) & (2) & (3) & (4) & (5) & (6) & (7) & (8) & (9) \\
1 & SCPMK0903-00601 & Konsep dasar algoritma; analisis permasalahan berbasis logika. & Modalitas: Daring/Luring. Metode: Contextual Teaching and Learning (CTL) dan Problem Based Learning. Penugasan: Tugas 1--2 studi kasus sederhana. & 1 x 2 x 100' tatap muka; 1 x 2 x 70' tugas terstruktur/mandiri. & Mahasiswa menjelaskan pentingnya algoritma serta memecahkan studi kasus dengan mengidentifikasi input, proses, dan output. & Keaktifan diskusi; ketepatan jawaban tugas. Bentuk: Tugas studi kasus. Mengacu rubrik RTM. & Ketepatan pengertian algoritma; identifikasi input-proses-output studi kasus. & 2 \\
2 & SCPMK0903-00601 & Konsep dasar version control; prinsip-prinsip Kanban. & Modalitas: Daring/Luring. Metode: CTL dan Problem Based Learning. Penugasan: analisis macam-macam version control; membuat Kanban board aktivitas sehari-hari. & 1 x 2 x 100' tatap muka; 1 x 2 x 70' tugas terstruktur/mandiri. & Mahasiswa menjelaskan pentingnya version control dan memvisualisasikan aktivitas dengan Kanban board. & Keaktifan diskusi; ketepatan jawaban tugas. Bentuk: Tugas studi kasus. Mengacu rubrik RTM. & Ketepatan konsep version control; kesesuaian Kanban board dengan prinsip visualisasi alur kerja. & 2 \\
3 & SCPMK0903-00602 & Tipe data, variabel, input-output, sequence, operator (aritmatika, gabungan, increment, decrement, relasional, logika, kondisional, bitwise, casting). & Modalitas: Daring/Luring. Metode: CTL dan Project Based Learning. Penugasan: Tugas 3 studi kasus tipe data dan operator. & 1 x 2 x 100' tatap muka; 1 x 2 x 70' tugas terstruktur/mandiri. & Mahasiswa mengidentifikasi dan mengimplementasikan konsep tipe data, variabel, input-output, sequence, dan operator pada studi kasus. & Keaktifan diskusi; ketepatan jawaban tugas. Bentuk: Tugas studi kasus. Mengacu rubrik RTM. & Ketepatan penjelasan tipe data, variabel, input-output, sequence, dan operator. & 2 \\
4 & SCPMK0903-00601, SCPMK0903-00602 & Kuis 1: materi minggu 1--3. & Ujian teori pilihan ganda dan/atau essay, closed book. & 1 x 2 x 100' ujian. & Mahasiswa mengerjakan soal kuis materi minggu 1--3 secara mandiri. & Ujian. Bentuk: Kuis 1. Mengacu rubrik RTM. & Ketepatan jawaban dengan kunci jawaban. & 10 \\
5 & SCPMK0704-00603 & Pemilihan 1 (if, if...else, if..else if..., switch..case); ekspresi logika. & Modalitas: Daring/Luring. Metode: CTL dan Project Based Learning. Penugasan: Tugas 4 studi kasus pemilihan 1 dengan flowchart. & 1 x 2 x 100' tatap muka; 1 x 2 x 70' tugas terstruktur/mandiri. & Mahasiswa menjelaskan konsep pemilihan 1 dan mengimplementasikannya dengan ekspresi logika untuk memecahkan studi kasus. & Keaktifan diskusi; ketepatan jawaban tugas. Bentuk: Tugas studi kasus. Mengacu rubrik RTM. & Ketepatan sintaks pemilihan 1; ketepatan penggambaran flowchart pemilihan. & 2 \\
6 & SCPMK0704-00603 & Pemilihan bersarang. & Modalitas: Daring/Luring. Metode: CTL dan Project Based Learning. Penugasan: Tugas 5 studi kasus pemilihan sederhana dan bersarang dengan flowchart. & 1 x 2 x 100' tatap muka; 1 x 2 x 70' tugas terstruktur/mandiri. & Mahasiswa mengimplementasikan algoritma pemilihan bersarang untuk menyelesaikan studi kasus. & Keaktifan diskusi; ketepatan jawaban tugas. Bentuk: Tugas studi kasus. Mengacu rubrik RTM. & Ketepatan konsep pemilihan bersarang; ketepatan flowchart studi kasus pemilihan bersarang. & 2 \\
7 & SCPMK0704-00603 & Konsep perulangan sederhana. & Modalitas: Daring/Luring. Metode: CTL dan Project Based Learning. Penugasan: Tugas 6 studi kasus perulangan bagian 1 dengan flowchart. & 1 x 2 x 100' tatap muka; 1 x 2 x 70' tugas terstruktur/mandiri. & Mahasiswa menjelaskan konsep perulangan bagian 1 dan mengimplementasikannya untuk memecahkan studi kasus. & Keaktifan diskusi; ketepatan jawaban tugas. Bentuk: Tugas studi kasus. Mengacu rubrik RTM. & Ketepatan konsep perulangan sederhana; ketepatan penggambaran flowchart perulangan. & 2 \\
8 & SCPMK0903-00601, SCPMK0903-00602, SCPMK0704-00603 & UTS: materi minggu 1--7. & Ujian teori pilihan ganda dan/atau essay, closed book. & 1 x 2 x 100' ujian. & Mahasiswa mengerjakan soal UTS materi minggu 1--7 secara mandiri. & Ujian. Bentuk: UTS. Mengacu rubrik RTM. & Ketepatan jawaban dengan kunci jawaban. & 30 \\
9 & SCPMK0704-00604 & Konsep array; array 1 dimensi; studi kasus searching dan sorting sederhana. & Modalitas: Daring/Luring. Metode: CTL dan Project Based Learning. Penugasan: Tugas 8 deklarasi dan operasi array 1 dimensi pada kasus searching. & 1 x 2 x 100' tatap muka; 1 x 2 x 70' tugas terstruktur/mandiri. & Mahasiswa menjelaskan konsep array 1 dimensi dan mengimplementasikannya untuk kasus searching dan sorting. & Keaktifan diskusi; ketepatan jawaban tugas. Bentuk: Tugas studi kasus. Mengacu rubrik RTM. & Ketepatan konsep array 1 dimensi; kebenaran solusi searching/sorting sederhana. & 1.5 \\
10 & SCPMK0704-00604 & Array 2 dimensi; studi kasus matriks. & Modalitas: Daring/Luring. Metode: CTL dan Project Based Learning. Penugasan: Tugas 9 studi kasus array 2 dimensi; penyimpanan nilai array 1 dan 2 dimensi pada proyek. & 1 x 2 x 100' tatap muka; 1 x 2 x 70' tugas terstruktur/mandiri. & Mahasiswa menjelaskan konsep array 2 dimensi dan mengimplementasikannya untuk menyelesaikan studi kasus matriks. & Keaktifan diskusi; ketepatan jawaban tugas. Bentuk: Tugas studi kasus. Mengacu rubrik RTM. & Ketepatan konsep array 2 dimensi; kebenaran solusi studi kasus matriks. & 1.5 \\
11 & SCPMK0704-00603 & Perulangan bersarang. & Modalitas: Daring/Luring. Metode: CTL dan Project Based Learning. Penugasan: Tugas 7 studi kasus perulangan bersarang dengan flowchart; identifikasi fitur perulangan bersarang pada proyek. & 1 x 2 x 100' tatap muka; 1 x 2 x 70' tugas terstruktur/mandiri. & Mahasiswa menjelaskan konsep perulangan bersarang dan mengimplementasikannya untuk memecahkan studi kasus/proyek. & Keaktifan diskusi; ketepatan jawaban tugas. Bentuk: Tugas studi kasus. Mengacu rubrik RTM. & Ketepatan konsep perulangan bersarang; ketepatan flowchart studi kasus perulangan bersarang. & 2 \\
12 & SCPMK0704-00603, SCPMK0704-00604 & Kuis 2: materi minggu 9--11. & Ujian teori pilihan ganda dan/atau essay, closed book. & 1 x 2 x 100' ujian. & Mahasiswa mengerjakan soal kuis materi minggu 9--11 secara mandiri. & Ujian. Bentuk: Kuis 2. Mengacu rubrik RTM. & Ketepatan jawaban dengan kunci jawaban. & 10 \\
13 & SCPMK0704-00604 & Konsep fungsi; fungsi iteratif; deklarasi dan pemanggilan fungsi. & Modalitas: Daring/Luring. Metode: CTL dan Project Based Learning. Penugasan: Tugas 10 studi kasus flowchart yang memuat fungsi. & 1 x 2 x 100' tatap muka; 1 x 2 x 70' tugas terstruktur/mandiri. & Mahasiswa menjelaskan konsep fungsi dan mengimplementasikan fungsi berparameter/return value pada studi kasus. & Keaktifan diskusi; ketepatan jawaban tugas. Bentuk: Tugas studi kasus. Mengacu rubrik RTM. & Ketepatan konsep fungsi; ketepatan flowchart yang menggunakan fungsi. & 1 \\
14 & SCPMK0704-00604 & Fungsi rekursif; perbedaan fungsi iteratif dan rekursif. & Modalitas: Daring/Luring. Metode: CTL dan Project Based Learning. Penugasan: Tugas 11 studi kasus fungsi rekursif dengan flowchart. & 1 x 2 x 100' tatap muka; 1 x 2 x 70' tugas terstruktur/mandiri. & Mahasiswa menjelaskan konsep fungsi rekursif dan mengimplementasikannya untuk menyelesaikan studi kasus. & Keaktifan diskusi; ketepatan jawaban tugas. Bentuk: Tugas studi kasus. Mengacu rubrik RTM. & Ketepatan konsep fungsi rekursif; ketepatan flowchart penyelesaian kasus rekursif. & 1 \\
15 & SCPMK0704-00604 & Studi kasus terintegrasi: sequence, pemilihan, perulangan, array, dan fungsi pada proyek. & Modalitas: Daring/Luring. Metode: CTL dan Project Based Learning. Penugasan: identifikasi fitur proyek yang memerlukan fungsi dan array, digambarkan dalam flowchart. & 1 x 2 x 100' tatap muka; 1 x 2 x 70' tugas terstruktur/mandiri. & Mahasiswa menyusun algoritma penyelesaian studi kasus terintegrasi dalam bentuk flowchart. & Keaktifan diskusi; ketepatan jawaban tugas. Bentuk: Tugas studi kasus. Mengacu rubrik RTM. & Kelengkapan dan kebenaran flowchart studi kasus terintegrasi. & 0.5 \\
16 & SCPMK0704-00604 & Responsi dan review studi kasus terintegrasi. & Modalitas: Daring/Luring. Metode: CTL, diskusi, dan review hasil tugas. & 1 x 2 x 100' tatap muka; 1 x 2 x 70' tugas terstruktur/mandiri. & Mahasiswa memperbaiki flowchart studi kasus terintegrasi berdasarkan umpan balik dosen. & Keaktifan diskusi; ketepatan jawaban tugas. Bentuk: Tugas studi kasus. Mengacu rubrik RTM. & Ketepatan perbaikan flowchart berdasarkan umpan balik. & 0.5 \\
17 & SCPMK0704-00603, SCPMK0704-00604 & UAS: materi minggu 1--16. & Ujian teori pilihan ganda dan/atau essay, closed book. & 1 x 2 x 100' ujian. & Mahasiswa mengerjakan soal UAS materi minggu 1--16 secara mandiri. & Ujian. Bentuk: UAS. Mengacu rubrik RTM. & Ketepatan jawaban dengan kunci jawaban. & 30 \\
\end{longtblr}
